\include{zLaTex-Preamble}


\newcommand{\bib}{

\begin{thebibliography}{9}

\bibitem{george}
  George, Andrew. 
  \emph{The Epic of Gilgamesh: The Babylonian Epic Poem and Other Texts in Akkadian and Sumerian}. 
  Penguin, 2002.
  Illustrated, reprint, reissue. 
  228 pages.

\bibitem{hoffner}
  Hoffner, Harry A., Jr.
  \emph{Hittite Myths}
  Writings from the Ancient World 2. 
  Atlanta: Scholars Press. 1990.

\bibitem{johnston}
  Johnston, S. I. 
  \emph{Religions of the Ancient World. A Guide.}
  Cambridge. Ma. 
  Harvard Univ. Press, 2004.

\bibitem{singer}
  Singer, Itamar.  
  \emph{Hittite prayers}.
  edited by Harry A. Hoffner, Jr.  
  ISBN 1-58983-032-6.

\bibitem{van_death}
 Bremer, J.M., van den Hout, Th.P.J., Peters, R.
 \emph{Death as a Privilege. The Hittite Royal Funerary Ritual} 
 Hidden Futures. Death and Immortality in Ancient Egypt, Anatolia, the Classical, Biblical and Arabic-Islamic World 
 Amsterdam. 37-75.
 1994.

\bibitem{van}
  van den Hout, Th PJ. 
  ``3. BIOGRAPHY AND AUTOBIOGRAPHY.'' 
  /emph{The Context of Scripture: Canonical compositions from the biblical world 1}
  1997.

\end{thebibliography}
}


\title{\papertitle}
\author{Joshua Miller}
\date{March 14, 2014}

\linespread{2}

\begin{document}

%% \maketitle
%% \newpage

~

As the first literate inhabitants of Mesopotamia\cite[p.20]{george},
the Sumerians provide an interesting point of comparison for
successive civilizations and the way their deities were represented
through the literature that they left behind.  To one extent, the
literature of the Hittites represents their gods in a similar manner
to those of the Sumerians, but to a greater extent, the Hittite gods
are presented with more pragmatism than their Sumerian counterparts.

Both Hittite and Sumerian mythologies present their gods as
anthropomorphic beings, with very human desires and emotions.
%% ~~~
A Sumerian example of a god exhibiting such human behaviors can be
found in the myth of Ishtar and the Bull of Heaven. The story begins
as the goddess Ishtar proposes marriage between to Gilgamesh: ``let me
be your bride and you shall be my husband'' \cite[p.102]{george}. To
this, Gilgamesh responds that he ``will not. How would it go with me?
Your lovers have found you like a brazier which smoulders in the cold,
a backdoor which keeps out neither squall of wind nor storm ... Which
of your lovers did you ever love for ever?''\cite[p.103]{george}.
%% ~~~
This direct insult incites Ishtar to seek a revenge so saturated with
emotion that it is impossible to deny the human nature of the goddess.
%% ~~~
She pleas for permission to release the Bull of Heaven unto Uruk, and
in turn kills at least three hundred innocent people. Ishtar's actions
can be characterized only as reactionary and not part of a divinely
just moral scheme.
%% ~~~
This passage alone demonstrates the extent to which the gods are
depicted almost as supernatural humans.
%% ~~~
In the Hittite disappearing deity myth, we see a similar human nature
in the gods.  Although many versions of this story forgo mention of
the cause, each one involves an enraged god's disappearance and the
parallel failure of agriculture, fertility, and water sources.
%% ~~~
In the Telipinu version, for example, in his anger ``Telipinu too went
away and removed grain, animal fecundity, luxuriance, growth, and
abundance''\cite[p.14]{hoffner} and the other Hittite gods became
worried. The world remained dessicate of resources until wax was
applied to a bee sting, soothing the delinquent god's anger.
%% ~~~
In both of the myths presented above, we see a shared anthropotheistic
representation of the deities between cultures. Moreover, these
stories serve to illustrate that both the Sumerian and Hittite deities
were representable in a physical form (the ability of the bee to sting
Telipinu, and the carnal request for a marriage to Isthar) and that
the pantheon was physitheistic. 
%% ~~~
Further evidence of this can be found in the mythos of the Sun God,
the Cow, and the Fisherman, as well as the tale of Appu and his
sons. The sun god not only attains a physical form when appearing to
Appu, ``as he changed himself into a young man, came to him, and
questioned him''\cite[p.83]{hoffner}, but again ``[became] a young
man, came down from the sky''\cite[p.85]{hoffner}. 
%% ~~~
As in the Sumerian myth involving Ishtar, we see that the gods are
capable of entering into corporeal encounters.  The shared
anthropopathy is hardly surprising as it was normative of religions
during the time period.

However, it is interesting to note that Hittite literature commonly
mentions the transformation of human king to a god whereas Sumerian
literature is absent of such rebirths. In Hittite texts, i.e. prayers,
edicts, etc., it was common to say that a king ``became
god''\cite[p.75,95,98] {singer}, \cite[p.196]{van}[etc.] upon his
death.  Furthermore, this acquired divinity may have been extended to
other members of the royal family.  Any questionability of actual
divine state of deceased kings, as opposed to being a simple formulaic
epithet devoid of substantial belief that kings became gods, is
quelled by the ``intricate royal funerary ritual lasting for two weeks
and heavily laden with symbolism'' \cite[p.38]{van_death} and the cult
worship of ``statues of deceased kings [that] had been erected in
temples and received offerings''\cite[p.45]{van_death}.

In stark contrast, every Sumerian mortal was consigned to post-death
existence in the morbid Netherworld, regardless of rank. 
%% ~~~
Although the legend of Gilgamesh curiously explains that ``Two thirds
they [the gods] made him [Gilgamesh], god and one third
man,'' \cite[p77]{george} his ascension to full divinity is
nevertheless an impossibility.
%% ~~~
After the death of his companion, Enkidu, Gilgamesh begins to fear
death and takes on a quest for immortality. Even ``the great king is
after all an ordinary mortal.''\cite[p.47]{george}. 
%% ~~~
While the Sumerian mythology presents gods in a separate sphere of
existence from humans, the Hittite literature blurs this distinction.



%% Hittite utility                

``There is no trace of any spontaneous prayer or of any prayer that
was purely adorative.'' \cite[p.359]{johnston}








\newpage
\bib
\end{document}



